\section{Dynamic chunking 的证据：何时学习边界胜过硬编码 token}

H-Net 的机制看起来合理，但“可学习”并不自动等于“更好”。本章用三组证据区分不同问题：第一组问 learned chunk 是否在足够数据后超过 BPE；第二组问即便输入已经是 BPE，outer stage 为什么仍偏好 Mamba；第三组把方法放到 DNA，观察没有成熟 tokenizer heuristic 时 scaling slope 是否改变。读图必须同时看 matched compute、crossover 与 domain boundary。

\subsection{Compute-matched crossover：先输后赢是 bitter lesson 的典型形状}

下图横轴是 total training bytes，纵轴是 bits per byte，模型 FLOPs 与大致规模被匹配。黑色 Transformer 直接使用 BPE；一层 byte H-Net 初期更差，因为它必须同时学语言与边界，数据足够后曲线交叉；两层 H-Net 进一步学习多级 abstraction，但也引入更多参数和优化难度。

\lecturefigure{slide-027.jpg}{Dynamic chunking scaling：isobitropic BPE、单层与多层 H-Net 的 compute-matched 比较}{Stanford 官方录像，00:40:40--00:44:09}

\begin{knowledgebox}{读图：三步检查 crossover}
第一，确认曲线是否在近似相同 FLOP budget 下比较，而不是只匹配参数；第二，观察 learned-chunk model 在低数据区的劣势，承认 hand-engineered BPE 提供了强 inductive bias；第三，关注高数据区 slope 与交叉点，判断额外 scaling 是否偿还了学习边界的成本。
\end{knowledgebox}

可把这类经验 scaling 写成
\[
\mathcal{L}(C)\approx \mathcal{L}_\infty + aC^{-\alpha},
\]
其中，$C$ 是训练 compute 或近似等价的数据预算，$\mathcal{L}_\infty$ 是不可约损失，$a$ 是常数，$\alpha$ 是 scaling exponent。图中的“scale better”不是说某个点更低，而是随着 $C$ 增大，learned hierarchy 的曲线保持更有利的下降趋势。有限区间拟合仍可能受 optimizer、data mixture 与 model sizing 影响。

\teachervoice{00:42:37--00:44:21，Gu 用 bitter lesson 解释 crossover：手工 feature 在 data-constrained regime 帮助很大；learned chunk 需要先花数据学会分段，但一旦学到比 BPE 更适合任务的单位，就可能继续受益于 scaling。}

\begin{warningbox}{“学习 token”并不保证无限外推}
当前曲线只覆盖有限规模和数据区间。更深 hierarchy 可能因 routing collapse、参数分配、优化不稳定或系统效率提前饱和。判断下一数量级是否继续占优，需要真正扩大模型和数据，而不是把直线延长到图外。
\end{warningbox}

\subsection{最关键 ablation：输入已是 BPE，outer Mamba 仍然重要}

如果 SSM 的优势只来自“byte sequence 太长”，那么把输入换成 BPE token 后，outer encoder 使用 Transformer 应当足够。图中却显示：即便整个系统操作 BPE，outer stages 加入 Mamba 仍明显改善 BPB。这里 outer stage 的任务不是简单读取输入，而是为 routing 制造可压缩表示；有限 state 的 inductive bias 因而可能直接帮助 abstraction formation。

\lecturefigure{slide-028.jpg}{在 BPE 输入上，H-Net outer stage 的 Mamba ablation 仍显示压缩式归纳偏置}{Stanford 官方录像，00:44:10--00:46:09}

\begin{importantbox}{本课最强的“compression is a feature”证据}
当输入单位已经是 BPE 时，resolution 解释不足以说明 outer Mamba 的收益。更合理的假设是：routing 前的 encoder 需要把局部历史整合成适合分块的表示，而 finite recurrent state 迫使网络形成 summary。这个结果仍是特定 H-Net ablation，却比单纯 byte benchmark 更直接地支持 compression-as-inductive-bias。
\end{importantbox}

\begin{warningbox}{Ablation 仍不能证明因果机制唯一}
Mamba 与 Transformer outer stage 可能还在 parameterization、normalization、optimization、kernel 或 effective depth 上不同。结果说明“使用 Mamba 的系统更好”，但不能单独证明全部增益只来自有限状态压缩。更严格的机制验证需要 matched architecture 与 representation probe。
\end{warningbox}

\subsection{dnaHNet：当没有语义 tokenizer 时，hierarchy 改变 scaling slope}

DNA 没有像英语空格、词根和词频那样成熟的 segmentation heuristic。图中每个点对应不同 FLOP budget，并在该 budget 下 sweep model/data allocation。H-Net 曲线不仅整体平移，还显示不同 slope；讲者据此认为 learned hierarchy 找到了 off-the-shelf tokenizer 无法提供的模式。

\lecturefigure{slide-029.jpg}{dnaHNet scaling laws：无成熟 tokenization heuristic 时 dynamic chunking 的增长趋势}{Stanford 官方录像，00:47:20--00:49:19}

本课按课堂日期采用 `dnaHNet` 的 arXiv v3（2026-04-09），因为它是 2026-04-16 前最新版本。图中引用的是 genomic autoregressive pretraining，不应把更低 loss 直接翻译为生物因果理解。真正的下游价值还取决于 regulatory element、variant effect 与 cell-type context 等任务验证。

\begin{knowledgebox}{为什么 DNA 是比 character English 更强的检验}
英语虽然可做 byte modeling，但已有成熟 BPE，研究者总可以问“为什么不用 tokenizer”。DNA 的单碱基序列缺少类似普适词边界，hand-engineered token 更难定义；若 learned chunking 在这里改变 scaling，更接近它最初要解决的原始数据问题。
\end{knowledgebox}

\teachervoice{00:47:17--00:49:18，讲者把 character English 称为仍带学术性质的长期方向，而把 DNA 视为 tokenizer 真正难以语义化的场景；这也是他认为 dynamic chunking 更有现实意义的原因。}

\subsection{本章小结}

H-Net 证据形成递进链：byte-level learned chunks 在足够数据后越过 BPE；BPE 输入上的 outer-stage ablation 仍偏好 Mamba，指向 compression inductive bias；DNA 在缺乏语义 tokenizer 时呈现更强 scaling 差异。链条支持“压缩可能帮助形成抽象”，但仍受有限规模、架构匹配与 domain transfer 限制。

